\documentclass{standalone}
\usepackage{circuitikz}
\usetikzlibrary{calc}
\definecolor{circuitnoteink}{HTML}{000000}
\definecolor{circuitnotemark}{HTML}{FE00FE}
\begin{document}
\begin{circuitikz}[american, line width=0.8pt]
\ctikzset{bipoles/length=1.2cm}
\foreach \x in {0,2,4,6} {\foreach \y in {0,-2,-4,-6} {\fill[gray, opacity=0.35] (\x,\y) circle (1.1pt);}}
\node[gray, font=\scriptsize] at (0,0.84) {1};
\node[gray, font=\scriptsize] at (2,0.84) {2};
\node[gray, font=\scriptsize] at (4,0.84) {3};
\node[gray, font=\scriptsize] at (6,0.84) {4};
\node[gray, font=\scriptsize, anchor=east] at (-0.84,0) {a};
\node[gray, font=\scriptsize, anchor=east] at (-0.84,-2) {b};
\node[gray, font=\scriptsize, anchor=east] at (-0.84,-4) {c};
\node[gray, font=\scriptsize, anchor=east] at (-0.84,-6) {d};
\coordinate (b2) at (2,-2);
\coordinate (b4) at (6,-2);
\coordinate (d2) at (2,-6);
\coordinate (d4) at (6,-6);
\draw (b2) to[R, l_=$R_{1}$, a^=$10\,\mathrm{k}\Omega$] (b4); % line 3
\draw (d2) to[R, l_=$R_{2}$, a^=$4.7\,\mathrm{k}\Omega$] (d4); % line 4
\draw[circuitnoteink] (0,0) -- (8,0); % line 6
\draw[circuitnoteink] (0,-4) -- (8,-4); % line 7
\draw[circuitnoteink] (0,-8) -- (8,-8); % line 8
\path (0,-1.593) rectangle (2.667,-0.407); % line 9
\node[anchor=west, inner sep=0, circuitnotemark, font=\large] at (0,-1) {X}; % line 9
\path (12,-6.921) rectangle (21.491,0.593); % line 10
\node[anchor=west, inner sep=0, circuitnotemark, font=\large] at (12,0) {X}; % line 10
\node[anchor=west, inner sep=0, circuitnotemark, font=\large] at (12,-0.487) {X}; % line 10
\node[anchor=west, inner sep=0, circuitnotemark, font=\large] at (12,-0.974) {X}; % line 10
\node[anchor=west, inner sep=0, circuitnotemark, font=\large] at (12,-1.46) {X}; % line 10
\node[anchor=west, inner sep=0, circuitnotemark, font=\large] at (12,-1.947) {X}; % line 10
\node[anchor=west, inner sep=0, circuitnotemark, font=\large] at (12,-2.434) {X}; % line 10
\node[anchor=west, inner sep=0, circuitnotemark, font=\large] at (12,-2.921) {X}; % line 10
\node[anchor=west, inner sep=0, circuitnotemark, font=\large] at (12,-3.408) {X}; % line 10
\node[anchor=west, inner sep=0, circuitnotemark, font=\large] at (12,-3.895) {X}; % line 10
\node[anchor=west, inner sep=0, circuitnotemark, font=\large] at (12,-4.381) {X}; % line 10
\node[anchor=west, inner sep=0, circuitnotemark, font=\large] at (12,-4.868) {X}; % line 10
\node[anchor=west, inner sep=0, circuitnotemark, font=\large] at (12,-5.355) {X}; % line 10
\node[anchor=west, inner sep=0, circuitnotemark, font=\large] at (12,-5.842) {X}; % line 10
\node[anchor=west, inner sep=0, circuitnotemark, font=\large] at (12,-6.329) {X}; % line 10
\node[anchor=south west, inner sep=2pt, circuitnotemark, font=\LARGE\bfseries] (circuittitle) at (current bounding box.north west) {X};
\path (circuittitle.south west) rectangle ++(3.658,0.853);
\end{circuitikz}
\end{document}
